\scanpage{45}
\mspage{-3.6-}
% archive sheet number: 44
for every $j\in\handwritten{\{}0,1,\ldots,n\handwritten{\}}$.

Note that $\handwritten{\lambda}_j$ is a section of $L_j\handwritten{\otimes}L_{\handwritten{\mu}(1)}^{-1}$ over the open set
$\handwritten{\unsure{\mathcal{U}|\mathscr{U}}}_R$. We now define
\[
\handwritten{\phi}:\handwritten{\mathcal{U}}_R\handwritten{\longrightarrow}\mathrm{PGL}(n)\mathbin{\mathrm{x}}\handwritten{\mathbb{A}}_R
\]
by its two components $p_1\circ\handwritten{\phi}$ and $p_2\circ\handwritten{\phi}$. The projective group
$\mathrm{PGL}(n)$, is, by definition,the affine open subset of
\[
\handwritten{\mathbb{P}}_{n^2+2n}=\operatorname{Proj}\handwritten{\mathbb{Z}}[a_{00},\ldots,a_{0n};\ a_{10},\ldots,a_{1n};\ \ldots;\ a_{n0},\ldots,a_{nn}]
\]
where $\det(a_{ij})\ne 0$. Therefore $p_1\circ\handwritten{\phi}$ is determined by
\[
\begin{aligned}
(p_1\circ\handwritten{\phi})^*(\underline{o}_{\handwritten{\mathbb{P}}}(1))&=L_{\handwritten{\mu}(1)}\\
(p_1\circ\handwritten{\phi})^*(a_{ij})&=X_i^{(j)}\handwritten{\otimes}(-1)^j\handwritten{\lambda}_j^{-1}.
\end{aligned}
\]
Finally the morphism $p_2\circ\handwritten{\phi}$ into affine space is determined by:
\[
(p_2\circ\handwritten{\psi})^*\left(x_i^{(n+k)}\right)=
\frac{D_{0,1,\ldots,\handwritten{\hat{i}},\ldots,n,n+k}}{D_{0,1,\ldots,\handwritten{\widehat{\mu(k)}},..,n,n+k}}\mathbin{\unsure{\cdot|\circ}}
\frac{\handwritten{\lambda}_i}{\handwritten{\lambda}_{\handwritten{\mu}(k)}}
\]
(It should be noted that the expression on the right is actually a
section of $\underline{o}_{\handwritten{\mathbb{A}_R}}$ ). Modulo the mechanical verifications that remain,
this proves the Proposition. \quad\uline{QED}

Of the morphisms constructed in the course of the above proof, it
is useful to \deleted{\unsure{set}} \handwritten{fix the notation}:
\[
\handwritten{\phi}_R=p_2\circ\handwritten{\phi}:\handwritten{\mathcal{U}}_R\handwritten{\longrightarrow}\handwritten{\mathbb{A}}^{n\unsure{\nu|v}-n}
\]

